\documentclass[11pt,a4paper]{article}

\usepackage[T1]{fontenc}
\usepackage[utf8]{inputenc}
\usepackage[margin=2.1cm]{geometry}
\usepackage{amsmath,amssymb,amsthm}
\usepackage{booktabs}
\usepackage{changepage}
\usepackage{enumitem}
\usepackage{xcolor}
\usepackage{listings}
\usepackage[colorlinks=true,linkcolor=black,citecolor=black,urlcolor=blue]{hyperref}

\definecolor{codebg}{RGB}{243,242,239}
\lstset{
  basicstyle=\ttfamily\footnotesize,
  backgroundcolor=\color{codebg},
  frame=none, breaklines=true, columns=fullflexible,
  xleftmargin=8pt, framexleftmargin=8pt, aboveskip=8pt, belowskip=8pt
}

\newtheorem{proposition}{Proposition}
\newtheorem{definition}{Definition}
\newtheorem{remark}{Remark}
\newtheorem{corollary}{Corollary}

\newcommand{\thc}{\theta_{\mathrm{crit}}}
\newcommand{\R}{\mathbb{R}}

\title{\textbf{Structural Identifiability of Model A}\\[0.3em]
\large Analysis, tooling, and the population-level result}
\author{Guido Putignano}
\date{\today}

\begin{document}
\maketitle

\begin{abstract}
\noindent
We analyse the structural identifiability of the scalar propionate model. Three results
follow. First, the model is \emph{globally structurally identifiable} from a single
patient under a generic (non-constant) trajectory: the obstruction is one of conditioning,
not of structure. Second, the identifiability degenerates sharply as
the trajectory approaches a steady state, and we give the rank condition and the growth of
the condition number. Third, and of most interest, at steady state a single patient is
\emph{not} identifiable while a cohort sharing population parameters \emph{is}, provided
the cohort contains specific anchoring genotypes. We state the anchor condition exactly.
The note also specifies the model in the input formats of three standard tools so that the
symbolic results can be verified mechanically.
\end{abstract}

\vspace{1em}

%=========================================================
\section{Setting}
%=========================================================

\subsection{The model}

\begin{equation}
\dot m = s(t) - \theta\, g(m) - k_e\, m,
\qquad g(m) = \frac{m}{K_m + m},
\qquad y = m,
\label{eq:A}
\end{equation}

with $p = (\theta, K_m, k_e) \in \R_{>0}^3$ the unknown parameter vector, $s(t)$ a known
input, and $y$ the measured output. The term $\theta g(m)$ is the Michaelis-Menten law
$v = V_{\max}[S]/(K_m + [S])$ with $V_{\max} = \theta$ and $[S] = m$, so $\theta$ is the
patient's residual maximal velocity, $K_m$ is the Michaelis constant of PCC, and $g$ is the
fractional saturation. The same form is called the Monod function, or the Hill function
with $n = 1$, elsewhere.

\subsection{What structural identifiability means}

\begin{definition}
Let $y(\cdot\,; p, s)$ denote the output of \eqref{eq:A} under parameter $p$ and input $s$.
The model is \emph{globally structurally identifiable} if, for generic $p$,
\[
y(\cdot\,; p_1, s) = y(\cdot\,; p_2, s) \ \text{ for all admissible } s
\quad \Longrightarrow \quad p_1 = p_2 ,
\]
\emph{locally identifiable} if the above holds within a neighbourhood of $p$, and
\emph{non-identifiable} otherwise.
\end{definition}

Three points determine how the result should be read.

\begin{enumerate}[leftmargin=*, itemsep=2pt]
\item It is a property of the \emph{model structure and the observation map}, not of any
      dataset. Noise, sampling frequency and cohort size do not enter.
\item It assumes the ideal case: continuous, exact, noise-free measurement of $y$ over an
      interval. It is therefore a \textbf{necessary condition}. A structurally
      non-identifiable parameter cannot be recovered by any amount or quality of data.
\item The converse fails. Structural identifiability says nothing about whether the
      problem is well conditioned in practice, and Section~\ref{sec:cond} shows that here
      it is not.
\end{enumerate}

%=========================================================
\section{Analytical result}
%=========================================================

\subsection{Global identifiability under a generic trajectory}

Since $y = m$ is measured exactly on an interval, $\dot m$ is known on that interval, and
$s$ is known. Hence the function
\begin{equation}
F(m) \;:=\; s(t) - \dot m(t) \;=\; \theta\, g(m) + k_e\, m
\label{eq:F}
\end{equation}
is known at every metabolite level $m$ visited by the trajectory. Clearing the denominator,
\begin{equation}
F(m) \;=\; \frac{k_e\, m^2 \;+\; (\theta + k_e K_m)\, m}{m + K_m} .
\label{eq:Frational}
\end{equation}

\begin{proposition}[Global structural identifiability]
\label{prop:global}
If the trajectory of \eqref{eq:A} visits an interval of positive length, then
$p = (\theta, K_m, k_e)$ is globally structurally identifiable.
\end{proposition}

\begin{proof}[Argument]
Knowing $F$ on an interval determines the rational function \eqref{eq:Frational}, hence its
coefficients in lowest terms: the denominator gives $K_m$, the leading numerator
coefficient gives $k_e$, and the remaining numerator coefficient gives $\theta + k_e K_m$,
from which $\theta$ follows. The map from $p$ to the coefficient triple
$(K_m,\, k_e,\, \theta + k_e K_m)$ is a bijection on $\R_{>0}^3$.
\end{proof}

\begin{remark}
This settles a question that is easily treated heuristically. The degeneracy of the
linearised regime, in which only one combination of $\theta$ and $k_e$ appears, is
\textbf{not} a structural obstruction. It is the leading-order behaviour of a
structurally identifiable problem evaluated near a degenerate operating point.
\end{remark}

\subsection{Rank condition}

Proposition~\ref{prop:global} requires a non-degenerate trajectory. The precise
requirement is a rank condition on the sensitivity matrix. Writing
\begin{equation}
\frac{\partial F}{\partial \theta} = g(m), \qquad
\frac{\partial F}{\partial K_m} = -\frac{\theta\, m}{(K_m + m)^2}, \qquad
\frac{\partial F}{\partial k_e} = m ,
\label{eq:sens}
\end{equation}
and collecting rows at distinct visited levels $m_1, \dots, m_r$ gives
$S \in \R^{r \times 3}$.

\begin{proposition}[Excitation requirement]
\label{prop:rank}
$\operatorname{rank} S = \min(r, 3)$ generically. Hence the trajectory must visit at least
\textbf{three distinct metabolite levels} for local identifiability of all three
parameters. One level gives rank $1$; two levels give rank $2$.
\end{proposition}

This is verified numerically in the accompanying code. It is the exact sense in which
``the patient must be excited'': not that a crisis is required, but that the metabolite
must occupy three distinguishable levels.

\subsection{Conditioning}
\label{sec:cond}

Rank is not the operative constraint; conditioning is. Table~\ref{tab:cond} reports the
condition number of the column-normalised sensitivity matrix as a function of the relative
excursion of $m$ about a nominal steady state, computed at
$(\theta, K_m, k_e) = (3000, 100, 5)$.

\begin{table}[htbp]
\centering
\small
\caption{Condition number against relative excursion of the metabolite.}
\label{tab:cond}
\begin{tabular}{@{}lll@{}}
\toprule
\textbf{Relative excursion} & \textbf{Levels visited} & \textbf{Condition number} \\
\midrule
$2\%$   & $9.8,\ 10,\ 10.2$   & $2.7 \times 10^{6}$ \\
$5\%$   & $9.5,\ 10,\ 10.5$   & $4.4 \times 10^{5}$ \\
$10\%$  & $9,\ 10,\ 11$       & $1.1 \times 10^{5}$ \\
$20\%$  & $8,\ 10,\ 12$       & $2.8 \times 10^{4}$ \\
$50\%$  & $5,\ 10,\ 15$       & $5.3 \times 10^{3}$ \\
$100\%$ & $10,\ 20,\ 30$      & $1.5 \times 10^{3}$ \\
$500\%$ & $10,\ 60,\ 110$     & $1.5 \times 10^{2}$ \\
\bottomrule
\end{tabular}
\end{table}

\begin{corollary}
The condition number scales approximately as the inverse square of the relative excursion.
A clinically well-controlled patient, whose metabolite varies by a few per cent between
visits, presents a problem conditioned at $10^{5}$ to $10^{6}$, which is beyond the
precision of the assay.
\end{corollary}

This is the correct statement of the difficulty. The problem is not that information is
structurally absent, but that its magnitude is inversely quadratic in the very excursion
that clinical management is designed to suppress.

%=========================================================
\section{The population result}
%=========================================================

The preceding analysis concerns a single patient with an exciting trajectory. The clinically
realistic case is the opposite: a cohort of patients each observed close to their own steady
state. We now show that identifiability can be recovered at the level of the cohort.

\subsection{Setting}

Consider $N$ patients indexed by $i$, sharing the population parameters $K_m$ and $k_e$,
each with an individual capacity $\theta_i$. Suppose patient $i$ is observed at a steady
state $(s_i, m_i)$, so that
\begin{equation}
s_i \;=\; \theta_i\, \frac{m_i}{K_m + m_i} \;+\; k_e\, m_i,
\qquad i = 1, \dots, N .
\label{eq:pop}
\end{equation}

This is a system of $N$ equations in $N + 2$ unknowns, namely $\theta_1, \dots, \theta_N$
together with $K_m$ and $k_e$. It is underdetermined with deficiency exactly $2$, regardless
of $N$: enlarging the cohort adds one equation and one unknown simultaneously.

\begin{remark}
This is the precise version of the informal claim that individual non-identifiability may be
resolved at population level. It is resolved only if two additional independent constraints
are supplied. Cohort size alone never suffices.
\end{remark}

\subsection{The anchors}

\begin{proposition}[Anchor condition]
\label{prop:anchor}
Suppose in addition that
\begin{enumerate}[label=(\roman*), leftmargin=*, itemsep=1pt]
\item the cohort contains at least one patient $j$ carrying a null genotype, so that
      $\theta_j = 0$; and
\item at least one patient $i$ is observed at two distinct steady states
      $(s_i, m_i)$ and $(s_i', m_i')$ with $m_i \neq m_i'$.
\end{enumerate}
Then $K_m$, $k_e$ and every $\theta_i$ are uniquely determined.
\end{proposition}

\begin{proof}[Argument]
From (i), \eqref{eq:pop} gives $s_j = k_e m_j$ directly, hence $k_e = s_j / m_j$: the null
genotype removes the enzymatic term entirely and measures the alternative channel in
isolation. With $k_e$ known, (ii) supplies two equations,
\[
s_i - k_e m_i = \theta_i \frac{m_i}{K_m + m_i}, \qquad
s_i' - k_e m_i' = \theta_i \frac{m_i'}{K_m + m_i'},
\]
in the two remaining unknowns $(\theta_i, K_m)$, which are generically solvable since
$m_i \neq m_i'$. With $K_m$ and $k_e$ known, each remaining equation in \eqref{eq:pop}
determines a single $\theta_i$.
\end{proof}

\subsection{Comments on the anchors}

\paragraph{Null genotypes are available.} In the Saudi series, the frameshift variant
\emph{PCCB} c.990dupT (p.E331Xfs) occurs in 7 families and is reported as a null allele.

\paragraph{The anchor does not depend on why the allele is null.} It is tempting to shore
up $\theta_j = 0$ with a molecular argument about what happens to the transcript, or to
worry about reports of residual propionate carboxylation in PCC-deficient cells and where
that activity comes from. Neither is needed, and the second is not a confound. Under
Proposition~\ref{prop:anchor} the requirement is not that no propionate is eliminated by
any route; it is only that whatever else eliminates it is not $\theta\, g(m)$. Any residual
elimination, of whatever molecular origin, is an alternative route and is absorbed into
$k_e$, which is exactly the quantity the null patients measure.

\paragraph{Carriers are not an anchor at all.} Obligate heterozygous parents are
asymptomatic, and it is natural to want to treat them as a known point on the activity
scale at roughly half of normal capacity. Two things stop this, and the first is decisive.

\emph{Measurement.} There is no assay that places a carrier on the scale. GeneReviews
states that quantitative plasma amino acids, urine organic acids, the acylcarnitine
profile and fibroblast enzymatic analyses are all unreliable for the detection of
heterozygotes. A value that cannot be measured in the individual cannot anchor anything,
and the half-of-normal figure is an inference from gene dosage rather than an observation.

\emph{Conditioning}, which would still bite if the first were solved. Being healthy,
carriers occupy a narrow range of $m$, over which total clearance is affine to within assay
error. Writing $\bar m$ for the centre of that range, the slope of the fit is
$\theta g'(\bar m) + k_e$, the same degenerate combination every stable patient contributes.
The intercept is $\theta\,(g(\bar m) - \bar m\, g'(\bar m)) = \theta\, g(\bar m)^2$, the
linear channel passing through the origin and contributing nothing to it, so a carrier does
in principle contribute a coefficient carrying $\theta$ by itself.

The conditioning of that contribution is what disqualifies them. Recovering $\theta$ from a
fitted intercept divides by $g(\bar m)^2$, and $g$ is increasing, so the amplification a
carrier suffers exceeds that of any patient by construction: carriers are the healthiest
people in the cohort and therefore the lowest on the saturation curve. The factor enters
squared, so a carrier whose $g(\bar m)$ is a third of a patient's is nine times worse placed
on that account alone, before the narrowness of their band enters through the extrapolation
to $m = 0$. Being asymptomatic is exactly what puts them there.

Carriers therefore fail twice over, and the two failures are independent. No assay places
them on the activity scale, and the one coefficient that would carry $\theta$ if an assay
existed is worse conditioned for them than for anyone the model is about. They constrain the
aggregate coefficient, not the separation, and they are not a substitute for condition (i).

\paragraph{Where the second anchor comes from.} Condition (ii) requires one patient observed
at two distinct steady states. Candidates present in routine care:
\begin{itemize}[leftmargin=*, itemsep=1pt]
\item a documented change in prescribed protein intake;
\item a change in carnitine dose, noting the confounding-by-indication caveat below;
\item weight-based dose recalculation following growth, which is quasi-exogenous;
\item introduction or interruption of cyclical intestinal antibiotics.
\end{itemize}

\paragraph{The weak point.} Both conditions require $s_i$ to be known. Substrate load is
only partially recorded: prescribed intake is documented, but gut flora contribution,
adherence and unrecorded catabolic episodes are not. Quantifying how uncertainty in $s$
propagates into $\theta_i$ is the first extension of this analysis.

\paragraph{Confounding by indication.} Carnitine dose is adjusted in response to the
patient's state, so treating it as an exogenous perturbation is invalid. Weight-based
recalculation and supply interruptions are the defensible instruments.

%=========================================================
\section{Verification by symbolic computation}
%=========================================================

Propositions~\ref{prop:global} and \ref{prop:rank} were obtained by hand and should be
confirmed mechanically. Three established tools decide structural identifiability for ODE
systems with rational right-hand sides, which \eqref{eq:A} is.

\begin{adjustwidth}{-1.0cm}{-1.0cm}
\begin{table}[htbp]
\centering
\footnotesize
\caption{Tools for structural identifiability of rational ODE models.}
\label{tab:tools}
\begin{tabular}{@{}p{0.17\textwidth} p{0.20\textwidth} p{0.28\textwidth} p{0.30\textwidth}@{}}
\toprule
\textbf{Tool} & \textbf{Platform} & \textbf{Method} & \textbf{Notes} \\
\midrule
\texttt{Structural\-Identifiability.jl}
& Julia, open source
& Differential algebra with Gr\"obner bases for global identifiability; Jacobian of
truncated power series for local
& Most modern of the three. Uniquely, it can also return \emph{identifiable functions of
the parameters}, i.e.\ which combinations are recoverable when individual parameters are
not. This is the function that should confirm the degenerate combination directly. Also
supports reparametrisation. \\
\addlinespace
\textsc{SIAN}
& Maple; a Julia port exists
& Hybrid symbolic--numeric. Randomised evaluation of Taylor coefficients with Gr\"obner
basis computation
& Returns local and global identifiability separately, with a user-specified probability of
correctness (e.g.\ $0.99$). Typically faster than pure differential algebra on
moderate systems. \\
\addlinespace
\textsc{DAISY}
& REDUCE
& Differential algebra: characteristic sets by Ritt pseudodivision, yielding
input--output relations whose coefficients form an exhaustive summary
& The original reference implementation. Conceptually the closest to the hand argument in
Section 2, since the input--output relation is precisely \eqref{eq:Frational} cleared of
denominators. Older and less convenient than the alternatives. \\
\bottomrule
\end{tabular}
\end{table}
\end{adjustwidth}

\begin{remark}[Why the \texttt{find\_identifiable\_functions} capability matters]
If the analysis of Section 2 is correct, then asking for identifiable functions under a
constant input should return the single combination $\theta/K_m + k_e$, whereas asking under
a generic time-varying input should return the three parameters individually. That contrast
is the mechanical confirmation of the entire argument, and it is a one-line computation.
\end{remark}

\subsection{Model specifications}

\paragraph{StructuralIdentifiability.jl.}
\begin{lstlisting}
using StructuralIdentifiability

ode = @ODEmodel(
    m'(t) = s(t) - theta * m(t) / (Km + m(t)) - ke * m(t),
    y(t)  = m(t)
)

assess_local_identifiability(ode)
assess_identifiability(ode)
find_identifiable_functions(ode)
\end{lstlisting}

\paragraph{SIAN, Maple.}
\begin{lstlisting}
read "IdentifiabilityODE.mpl":
sigma := [
  diff(m(t), t) = s(t) - theta*m(t)/(Km + m(t)) - ke*m(t),
  y(t) = m(t)
]:
IdentifiabilityODE(sigma, GetParameters(sigma), p_mod = 0.99);
\end{lstlisting}

\paragraph{DAISY.} The input--output relation obtained by clearing denominators in
\eqref{eq:A} is
\begin{equation}
(K_m + y)\,\dot y \;=\; (K_m + y)\, s \;-\; \theta\, y \;-\; k_e\, y\,(K_m + y),
\label{eq:io}
\end{equation}
which expands to a polynomial relation in $y, \dot y, s$ whose coefficients are
\[
\{\, K_m,\ \ \theta + k_e K_m,\ \ k_e \,\},
\]
the exhaustive summary. Injectivity of $p \mapsto$ (this triple) is Proposition
\ref{prop:global}. The DAISY input file in the accompanying archive encodes
\eqref{eq:A} directly; the syntax should be checked against the current DAISY manual, as
the format has changed between releases.

\subsection{Cases to run}

\begin{enumerate}[leftmargin=*, itemsep=2pt]
\item \textbf{A1.} Baseline, generic input $s(t)$. Expected: all three globally identifiable.
\item \textbf{A2.} Constant input, $s(t) \equiv s_0$. Expected: only the degenerate
      combination identifiable. This is the formal statement of the clinical situation.
\item \textbf{A3.} Two patients sharing $(K_m, k_e)$ with individual $\theta_1, \theta_2$,
      constant inputs. Tests whether cohort structure alone helps, which
      Section 3 predicts it does not.
\item \textbf{A4.} As A3, with $\theta_2 = 0$ imposed (null anchor). Expected: $k_e$ becomes
      identifiable.
\item \textbf{A5.} Carnitine as an input on the alternative channel,
      $k_e \to k_e + \kappa\, u_c(t)$.
\item \textbf{A6.} Therapy as an input on the enzymatic channel,
      $\theta \to \theta + \Delta(t)$.
\item \textbf{B.} Model B, with $k_e m$ replaced by $k_e m/(K_e + m)$. Expected: weaker,
      since the two channels acquire the same functional form.
\item \textbf{C.} Model C, the two-state carnitine depletion system.
\end{enumerate}

Cases A2 and A4 are the ones that matter. A2 should reproduce the degeneracy formally, and
A4 should show that a null genotype restores identifiability of the shared parameter.

%=========================================================
\section{Open points}
%=========================================================

\begin{enumerate}[leftmargin=*, itemsep=3pt]
\item \textbf{Propagation of uncertainty in $s$.} Both anchors assume $s_i$ known. It is
      not. The dependence of the recovered $\theta_i$ on misspecification of $s$ is the
      first quantity to bound.

\item \textbf{Optimal experimental design under a state constraint.} Given that the
      condition number scales as the inverse square of the excursion, and that the excursion
      is bounded above by safety, the design question is how much information can be
      extracted subject to never approaching $\thc$. This is where the analysis meets
      dual control.

\item \textbf{Identifiability of the decision rather than the parameters.} The clinical
      quantity is the sign of $\theta - \thc$. Whether that comparison is better conditioned
      than either parameter separately is open; the opposite signs with which $\theta$ and
      $k_e$ enter the margin suggest it may be worse.

\item \textbf{Multi-parameter estimator design.} The scalar Immersion and Invariance
      construction requires a sign condition on a scalar. With three parameters entering
      simultaneously the attractivity requirement becomes a matrix condition, and the
      existence of a suitable design function is not guaranteed.
\end{enumerate}

\end{document}
